% TEAM: check every statement here against Section 4 once it is written.
\section{Discussion}\label{sec:discussion}

\subsection{Why the CIDA baseline underperforms}\label{sec:disc-baseline}
The baseline stays below the paper on every task (\cref{sec:exp-vanilla}).
All of these losses come from one mechanism.

\paragraph{The model trains on its own mistakes.}
Only one image per private class is labelled, so the method labels the rest of the target data itself.
Each target image takes the class of its nearest guide, and $\mathcal{L}_{a1}$ then pulls the image towards that guide.
When an early pseudo-label is wrong, $\mathcal{L}_{a1}$ moves the feature towards the wrong guide.
The next pseudo-labels are then worse, and the error grows with every epoch.
This effect is known as confirmation bias in pseudo-labelling~\citep{arazo2020pseudo}.
Without target labels, such self-training can also collapse onto a few classes~\citep{liang2020shot}, which is what happens on VisDA.
This is why the best epoch of every baseline task lies between epoch 0 and epoch 12, and accuracy falls after it.
Four details of the method make the first pseudo-labels wrong or let the errors spread.

\paragraph{1. Source BatchNorm statistics.}
$f_t$ starts as a copy of $f_s$, so it normalises target images with statistics computed on source images.
These statistics carry domain information, so they do not fit the target domain~\citep{li2018adabn}.
The private guides come from real target images, while the shared guides come from source class means.
Shared target images therefore lie closer to a private guide than to their own shared guide.
On Digits M$\to$U, 74 to 78\% of the shared samples receive private pseudo-labels for this reason.
With target statistics this share falls to 14 to 21\%.

\paragraph{2. Projections that are not the identity.}
The paper pretrains $f_e$ and $f_d$ as plain MLPs towards the identity, but an error remains.
After 2{,}000 steps the relative error is about 13\% on Digits and 24\% on VisDA.
The last layer also carries an activation, which distorts the negative values of BatchNorm features.
The shared guides $f_e(\mu_s)$ are therefore misplaced before adaptation starts.
On Digits this lowers shared accuracy to between 28 and 85\% across seeds before the first adaptation step.
A residual branch whose last layer starts at zero would make the projection exactly the identity~\citep{zhang2019fixup}, which the paper does not do.

\paragraph{3. $\mathcal{L}_{a1}$ on every target sample.}
As written, $\mathcal{L}_{a1}$ trains $f_t$ on the pseudo-label of every target sample, including the many wrong ones.
In the first epoch of the baseline, only 37.5 to 43.9\% of these labels are correct on Digits.
$\mathcal{L}_{a1}$ is the loss that turns wrong labels into moved features, so it drives the loop above.
Self-training methods for domain adaptation limit this by training only on confident pseudo-labels~\citep{zou2018cbst}.
With the other changes in place, shared accuracy on Digits M$\to$U at epoch 30 is 52.8 when $\mathcal{L}_{a1}$ uses every sample and 83.4 when it uses only $\Bt$.

\paragraph{4. Fine-tuning the whole ResNet-50.}
On Office-31 and VisDA the baseline trains the ResNet-50 together with the rest of the network.
On Office-31 the source domain holds about 900 images of 10 classes, and the network fits them to 100\% source accuracy.
In doing so it discards the ImageNet features that separate the unseen private classes.
Full fine-tuning is known to distort pretrained features and to lower accuracy under distribution shift~\citep{kumar2022finetuning}.
On A$\to$D, nearest-guide accuracy on the private classes is 72.4 with ImageNet features and only 27.0 to 33.1 after fine-tuning.
With weaker features, the noisy labels of $\mathcal{L}_{a1}$ reshape the features even further.

\paragraph{Our reproduction or the method.}
The paper leaves several settings open, such as epoch counts, so part of the gap could come from our choices in \cref{tab:choices}.
Two results suggest that the method as written is the main cause.
First, stopping earlier does not close the gap.
Even at the best epoch, chosen with target labels, ALL stays 1 to 19 points below the paper on Office-31 and 24 to 33 points below on VisDA.
Second, removing the four causes above recovers most of the gap (\cref{sec:disc-improvements}).

\subsection{How the improvements help}\label{sec:disc-improvements}

\subsection{Limitations}\label{sec:disc-limitations}
\begin{itemize}
  \item \textbf{Accuracy on the new classes.} Accuracy on the target-private classes (PRIV) remains below the values reported in the paper (\cref{sec:exp-improvements}).
  \item \textbf{Departures from the paper.} Our four changes alter the method. We report the method as written separately in \cref{sec:exp-vanilla}, so readers can compare the two versions.
  \item \textbf{Parts not implemented.} We did not implement Ours-d, the paper's variant that uses real source features during adaptation. We also did not implement the two-step upgrade from Amazon to DSLR and then to Webcam (A$\to$D$\to$W).
\end{itemize}
